\documentclass[aps,pra,reprint,amsmath,amssymb,superscriptaddress,longbibliography]{revtex4-2}

\usepackage{booktabs}
\usepackage{bm}
\usepackage{hyperref}
\usepackage{microtype}
\usepackage{xcolor}
\setcounter{MaxMatrixCols}{12}

\newcommand{\ii}{\mathrm{i}}
\newcommand{\ee}{\mathrm{e}}
\newcommand{\per}{\operatorname{per}}
\newcommand{\Tr}{\operatorname{Tr}}
\newcommand{\btheta}{\bm{\theta}}
\newcommand{\be}{\bm e}
\newcommand{\Ffour}{F_4}

\begin{document}

\title{Dimension-saturating local tomography from Fourier-cat dark events}

\author{Hainan Zhao}
\email{hainzhao@gmail.com}
\affiliation{Independent Researcher}

\date{\today}

\begin{abstract}
An exact dark event is a sensitive null test for a linear-optical device, but
its unprobed probability is quadratic in a small coherent error and has no
regular local inverse.  We show that positive and negative calibrated output
probes turn dark-event counts into linear phase-sensitive observables.  A
general signed-probe theorem gives their limiting Jacobian and proves that
one dark event contributes at most two real rows.  Number-state inputs face
an exact obstruction: in an $m$-mode Fourier interferometer, an
$(m-1)$-dimensional subspace of off-diagonal output generators is
indistinguishable from input phases, even after an arbitrary known output
analyzer.  A coherent Fourier-cat input removes this continuous gauge.
For every $m\geq2$ and every $n>2$ divisible by $m$, we construct
$m(m-1)/2$ dark outcomes and two fixed probes whose $m(m-1)$ signed
contrasts identify every off-diagonal Hermitian error locally.  Within the
weak-probe, regular-local protocol defined here, both the outcome and
scalar-contrast counts attain their dimension lower bounds; four programmed
settings suffice for every $m$, although we do not prove that four is the
minimum possible number of settings.  At $m=n=4$ we additionally
give a sparse exact implementation with Jacobian determinant
$-177147/256$.  We prove an $O(\epsilon^2)$ finite-probe bias bound and show
how internal coherent errors propagate to the estimated output generator.
The Fourier-cat suppression rule and displaced-null principle are known;
the contribution is their explicit bosonic identifiability construction
together with the Fock gauge theorem and exact dimension-saturating
certificates.
\end{abstract}

\maketitle

\section{Introduction}
\label{sec:introduction}

Destructive multiphoton interference produces output events with exactly
zero ideal probability.  Fourier suppression laws generalize the
Hong--Ou--Mandel effect to multiport devices
\cite{Tichy2010,Crespi2016,Dittel2018PRL,Dittel2018PRA,Bezerra2023}, and
forbidden-event populations have become practical witnesses of interference
and device quality \cite{Crespi2016,Wang2023,Sanz2026}.  A scalar
forbidden-event fraction, however, does not specify which coherent device
error produced the counts.

There is a more basic local difficulty.  If an event is dark at a nominal
device, its probability has zero derivative with respect to every coherent
error coordinate.  The leading quadratic response cannot distinguish
opposite errors.  The situation parallels the non-identifiability of exact
null measurements in quantum statistics.  Girotti, Godley, and Guţă showed
that a deliberate displacement restores identifiability and can support
optimal estimation in general pure-state models \cite{Girotti2024}.  Here we
develop a finite-mode, photon-counting implementation of that principle:
a known signed optical probe supplies a reference amplitude against which
unknown leakage interferes.

The central issue is not merely sensitivity but \emph{identifiability}.
Number-state preparation and number-resolved detection possess independent
input and output phase gauges, familiar from linear-optical
characterization \cite{RahimiKeshari2013}.  At a Fourier device, the
transported input-phase gauge lies entirely inside the nominally
off-diagonal output-error coordinates.  We prove that number-state
experiments cannot see these directions even if an arbitrary calibrated
analyzer is appended.  An exact four-photon search then saturates, rather
than evades, this rank bound.

The missing resource is an input phase reference.  We use equal
superpositions of all-bunched number states.  Fourier multiports and
multimode cat states have a substantial prior literature
\cite{VourdasDunningham2005}; the associated modular dark sectors also
follow from the general symmetry suppression framework
\cite{Dittel2018PRA}.  We use this known selection rule to produce a
dimension-saturating frame of simultaneously dark outcomes for arbitrary
mode number.  General near-identity qudit-unitary estimation has also been studied
with a controlled-qudit architecture \cite{EscandonMonardes2024}; our
construction instead uses bosonic Fourier interference, number-resolved
dark outcomes, and two fixed global probes.

Our contributions are:
\begin{enumerate}
\item a general count-level signed-probe theorem, including the two-real-row
bound per dark event and a rigorous $O(\epsilon^2)$ finite-angle bias bound;
\item an exact $(m-1)$-dimensional number-state gauge no-go theorem;
\item an outcome- and contrast-minimal rank-nine $F_4$ certificate, within
the same regular-local framework, using one four-photon Fock input;
\item for all $m\geq2$ and $n>2$ divisible by $m$, a
dimension-saturating Fourier-cat construction identifying all $m(m-1)$
off-diagonal Hermitian coordinates from $m(m-1)/2$ dark outcomes and four
programmed signed settings;
\item a sparse exact $F_4$ implementation with a rational full-rank
certificate;
\item a first-order adjoint rule connecting arbitrary internal coherent
circuit errors to the reconstructed output generator.
\end{enumerate}
The result is local tomography on the declared $m(m-1)$-dimensional
off-diagonal slice.  It is not global unitary tomography and does not
identify the $m-1$ output-diagonal $SU(m)$ generators.

\section{Model and parameter target}
\label{sec:model}

Let $U_0$ be an $m$-mode passive interferometer and
$\widehat U_0$ its action on a fixed $n$-photon sector.  For number
occupations $\bm r,\bm s$, the normalized transition amplitude is
\begin{equation}
 {\cal A}_{\bm r,\bm s}(U)
 =\frac{\per U[\bm s,\bm r]}
 {\sqrt{\prod_jr_j!\prod_ks_k!}}.
\label{eq:fock-amplitude}
\end{equation}
Rows and columns are repeated according to $\bm s$ and $\bm r$.

The unknown coherent output generator belongs to
\begin{align}
 {\cal H}_{\rm off}(m)
 &=\{H=H^\dagger:H_{aa}=0\},\nonumber\\
 \dim_{\mathbb R}{\cal H}_{\rm off}(m)&=m(m-1).
\label{eq:hoff}
\end{align}

For the explicit four-mode certificates we use
\begin{equation}
 (\Ffour)_{kj}=\frac{\ii^{kj}}2,\qquad 0\leq k,j<4.
\label{eq:f4}
\end{equation}
The twelve off-diagonal Hermitian generators are
\begin{align}
 X_{pq}&=|p\rangle\langle q|+|q\rangle\langle p|,\nonumber\\
 Y_{pq}&=-\ii|p\rangle\langle q|+\ii|q\rangle\langle p|,
 \qquad 0\leq p<q<4.
\label{eq:xy}
\end{align}
We fix the coordinate order
\begin{equation}
\btheta=(x_{01},y_{01},x_{02},y_{02},x_{03},y_{03},
x_{12},y_{12},x_{13},y_{13},x_{23},y_{23})
\label{eq:coordinate-order}
\end{equation}
and write $G(\btheta)=\sum_\mu\theta_\mu G_\mu$.  The unknown device is
locally
\begin{equation}
 U(\btheta)=\ee^{\ii G(\btheta)}\Ffour.
\label{eq:unknown-device}
\end{equation}
The normalization $\Tr X_{pq}^2=\Tr Y_{pq}^2=2$ fixes the meaning of the
reported coordinates and condition numbers.

The three traceless output-diagonal generators are deliberately excluded.
For an output occupation $\bm s$, a diagonal output unitary multiplies its
amplitude by a phase, leaving its number probability unchanged.  Moreover,
at a dark event
\begin{equation}
 \langle\bm s|\ii\widehat Z\,\widehat U_0|\psi\rangle
 =\ii z(\bm s)\langle\bm s|\widehat U_0|\psi\rangle=0 .
\label{eq:output-diagonal-zero}
\end{equation}
They therefore do not occur in the leading dark-event Jacobian.  A finite
noncommuting analyzer can create higher-order sensitivity to them, but that
is a different parameter model.

\section{Signed-probe identifiability}
\label{sec:signed-probe}

Consider events $e=(|\psi_e\rangle,|\bm s_e\rangle)$ that are dark at
$U_0$:
\begin{equation}
 \langle\bm s_e|\widehat U_0|\psi_e\rangle=0.
\label{eq:dark}
\end{equation}
For a real generator coordinate $\bm u$, define the complex amplitude
functional
\begin{equation}
 \ell_e\bm u
 =\langle\bm s_e|\ii\widehat G(\bm u)\widehat U_0|\psi_e\rangle .
\label{eq:ell}
\end{equation}

\subsection{Why an undisplaced null is insufficient}

Without a probe,
\begin{equation}
 p_e(\btheta)=|\ell_e\btheta|^2+O(\|\btheta\|^3).
\label{eq:quadratic-null}
\end{equation}
Thus $Dp_e(0)=0$.  More strongly, no collection of such unprobed dark
probabilities has a locally Lipschitz left inverse on an open parameter
neighborhood: along $\btheta=t\bm u$, a Lipschitz inverse would imply
$|t|\|\bm u\|=O(t^2)$, a contradiction.  This does not exclude a
nonregular higher-order inverse on a restricted domain; it rules out the
regular local tomography needed for a stable signed estimate.

\subsection{Calibrated central contrast}

Append a known probe $H(\bm h)$ after the unknown error and program both
signs:
\begin{equation}
 p_{e,\bm h}^{\sigma}(\btheta,\epsilon)
 =\left|
 \langle\bm s_e|
 \ee^{\ii\sigma\epsilon\widehat H(\bm h)}
 \ee^{\ii\widehat G(\btheta)}
 \widehat U_0|\psi_e\rangle
 \right|^2,\qquad \sigma=\pm1 .
\label{eq:probe-probability}
\end{equation}
We use the baseline-subtracted response
\begin{align}
 R_{e,\bm h}(\btheta,\epsilon)
 &=
 \epsilon^{-1}\big[
 p^+_{e,\bm h}(\btheta,\epsilon)
 -p^-_{e,\bm h}(\btheta,\epsilon)\nonumber\\
 &\hspace{2.2cm}
 -p^+_{e,\bm h}(0,\epsilon)
 +p^-_{e,\bm h}(0,\epsilon)\big].
\label{eq:response}
\end{align}
This is four times the convention
$[p^+-p^-]/(4\epsilon)$ often used for a central difference.

\paragraph*{Theorem 1 (signed-probe Jacobian).}
For every finite-dimensional photon sector,
\begin{equation}
 \left.
 \frac{\partial}{\partial\theta_\mu}
 \lim_{\epsilon\to0}R_{e,\bm h}(\btheta,\epsilon)
 \right|_{\btheta=0}
 =
 4\,\operatorname{Re}
 \left[(\ell_e\bm h)^*(\ell_e\be_\mu)\right].
\label{eq:signed-probe-theorem}
\end{equation}
At finite angle the derivative differs from
Eq.~\eqref{eq:signed-probe-theorem} by $O(\epsilon^2)$.

A proof and a norm-explicit remainder bound are given in
Appendix~\ref{app:signed-proof}.  Equation~\eqref{eq:signed-probe-theorem}
is a statement about observable probability differences, not an assumption
that complex transition amplitudes can be measured directly.  The probe
amplitude $\ell_e\bm h$ acts as a phase reference.

\paragraph*{Corollary 1 (two-row event bound).}
For one dark event, all limiting signed-probe rows lie in
\begin{equation}
 \operatorname{span}_{\mathbb R}
 \{\operatorname{Re}\ell_e,\operatorname{Im}\ell_e\}.
\label{eq:two-row-span}
\end{equation}
Hence one event contributes at most two independent real first-order
measurements, no matter how many probes are applied.

If a limiting $q\times d$ response matrix has full column rank, some
$d\times d$ row minor is nonsingular.  Continuity preserves that minor for
all sufficiently small nonzero calibrated $\epsilon$.  The ordinary local
inverse-function theorem applied to the selected $d$ responses then gives a
local differentiable left inverse.  This argument must not be confused with
the algebraic Jacobian conjecture, which concerns a polynomial self-map with
constant nonzero determinant and a \emph{global polynomial} inverse.  Nor
does full rank at one point imply global injectivity.  Our map need not be
polynomial, the determinant is evaluated locally, and only local
identifiability is claimed.

\section{The number-state gauge and its optimal rank}
\label{sec:fock-gauge}

\paragraph*{Theorem 2 (transported input-phase gauge).}
Let $F_m$ be the $m$-mode Fourier matrix.  For every real traceless diagonal
$D$, the generator
\begin{equation}
 K_D=F_mDF_m^\dagger
\label{eq:transported-gauge}
\end{equation}
has zero diagonal and is therefore an off-diagonal Hermitian generator.
Nevertheless, for every input number state $|\bm r\rangle$,
\begin{equation}
 \ee^{\ii\widehat K_D}\widehat F_m|\bm r\rangle
 =\widehat F_m\ee^{\ii\widehat D}|\bm r\rangle
 =\ee^{\ii\sum_jd_jr_j}\widehat F_m|\bm r\rangle .
\label{eq:gauge-no-go}
\end{equation}
Consequently $K_D$ is invisible to all output number probabilities, even
after an arbitrary known output analyzer.  The invisible subspace has
dimension $m-1$.

The proof is included in Appendix~\ref{app:fock-certificate}.  At $m=4$,
an explicit basis inside the coordinate order
\eqref{eq:coordinate-order} is
\begin{align}
 K_A&=X_{01}+X_{03}+X_{12}+X_{23},\nonumber\\
 K_B&=Y_{01}-Y_{03}+Y_{12}+Y_{23},\nonumber\\
 K_C&=X_{02}+X_{13}.
\label{eq:gauge-basis}
\end{align}
Thus Fock-input experiments can identify at most nine of the twelve
off-diagonal coordinates.

\subsection{Exact four-photon saturation}

The bound is tight.  Use input $\bm r=(1,1,1,1)$ and the nine
output--probe contrasts in Table~\ref{tab:fock-design}.  All five outputs
are exactly dark by the cyclic Fourier selection rule
\cite{Tichy2010,Dittel2018PRA}.
\begin{table}[t]
\caption{An exact rank-nine design.  Multiple probes on one output give
separate scalar contrasts.}
\label{tab:fock-design}
\begin{ruledtabular}
\begin{tabular}{cc}
output $\bm s$ & probes\\
\hline
$(0,0,1,3)$ & $X_{03},Y_{03}$\\
$(0,0,2,2)$ & $X_{02},Y_{02}$\\
$(0,0,3,1)$ & $X_{12},Y_{12}$\\
$(0,1,1,2)$ & $X_{23},Y_{23}$\\
$(0,2,2,0)$ & $Y_{02}$
\end{tabular}
\end{ruledtabular}
\end{table}

The normalized $9\times12$ limiting matrix is displayed and derived in
Appendix~\ref{app:fock-certificate}.  It annihilates exactly the three
vectors in Eq.~\eqref{eq:gauge-basis}; a $9\times9$ minor has determinant
$81/8$.  Therefore its rank is nine.  Nine scalar contrasts are necessary
for a nine-dimensional quotient, and five events are necessary by
Corollary~1.  The design attains both lower bounds.

\section{A coherent phase reference}
\label{sec:cat}

\subsection{Known Fourier-cat dark sectors}

Let
\begin{equation}
 F_m(k,j)=m^{-1/2}\omega^{kj},\qquad
 \omega=\ee^{2\pi\ii/m},
\end{equation}
and define the phase-twisted all-bunched cat
\begin{equation}
 |\Psi_\ell^{(m,n)}\rangle
 =\frac1{\sqrt m}\sum_{j=0}^{m-1}
 \omega^{-\ell j}|n\be_j\rangle .
\label{eq:general-cat}
\end{equation}
For an output occupation $\bm s$ of total size $n$, put
$Q(\bm s)=\sum_kks_k\pmod m$.  Direct evaluation gives
\begin{equation}
 {\cal A}_{\Psi_\ell,\bm s}(F_m)
 =
 \sqrt{\frac{n!}{\prod_ks_k!}}\,
 m^{(1-n)/2}\,
 \delta_{Q(\bm s),\ell}.
\label{eq:cat-sector}
\end{equation}
Thus all outputs outside one modular sector are dark.  The Fourier
multiport--cat connection is known from
Vourdas and Dunningham \cite{VourdasDunningham2005}, and
Eq.~\eqref{eq:cat-sector} is also an immediate cyclic-symmetry selection
rule \cite{Dittel2018PRA}.  We recall it to specify our phase convention and
resource; we do not claim it as a new suppression law.  A self-contained
derivation is given in Appendix~\ref{app:cat-certificate}.

\subsection{An all-mode dimension-saturating frame}

We now specialize to the charge-zero cat, take $n>2$ divisible by $m$, and
write
\begin{equation}
 \alpha_{\bm s}=
 \sqrt{\frac{n!}{\prod_as_a!}}\,m^{(1-n)/2}.
\end{equation}
For an output with $Q(\bm s)=c\neq0$, define the phase-free dark-amplitude
functional
\begin{equation}
 L_{\bm s}(H)=\alpha_{\bm s}\sum_{a=0}^{m-1}s_aH_{a,a-c},
\label{eq:general-dark-functional}
\end{equation}
where indices are modulo $m$.  Its physical amplitude derivative is
$\ell_{\bm s}(H)=\ii L_{\bm s}(H)$.

\paragraph*{Theorem 3 (all-mode Fourier-cat frame).}
For every $m\geq2$ and every $n>2$ divisible by $m$, there are
$m(m-1)/2$ dark outcomes and two fixed Hermitian probes $R,T$ for which
the two limiting signed contrasts per outcome have full rank on
${\cal H}_{\rm off}(m)$.  Thus four settings,
$\pm\epsilon R$ and $\pm\epsilon T$, suffice for local identification at
all sufficiently small calibrated nonzero $\epsilon$.

Here is the explicit construction.  Choose the cyclic separations
$c=1,\ldots,\lfloor m/2\rfloor$.  If $2c\not\equiv0\pmod m$, use all $m$
outcomes
\begin{equation}
 \bm s_{p,c}=(n-1)\be_p+\be_{p+c},\qquad 0\leq p<m.
\label{eq:selected-general-outcomes}
\end{equation}
If $m$ is even and $c=m/2$, retain only $0\leq p<m/2$.  Their charge is
$np+c\equiv c$, so all are dark.  Their total number is
$m(m-1)/2$.

Let $R$ be the real complete-graph probe,
\begin{equation}
 R_{ab}=1\ (a\neq b),\qquad R_{aa}=0.
\end{equation}
Orient every edge of cyclic separation $c<m/2$ by
\begin{equation}
 T_{a,a-c}=\ii,\qquad T_{a-c,a}=-\ii.
\label{eq:oriented-probe}
\end{equation}
For even $m$, orient the remaining antipodal edges as
$T_{p,p+m/2}=\ii$ for $0\leq p<m/2$, with the reverse entries fixed by
Hermiticity.  Then
\begin{align}
 L_{p,c}(R)&=\alpha n,&L_{p,c}(T)&=\ii\alpha n
 &&(2c\not\equiv0),\nonumber\\
 L_{p,c}(R)&=\alpha n,&L_{p,c}(T)&=\ii\alpha(n-2)
 &&(2c\equiv0),
\label{eq:general-reference-values}
\end{align}
where the selected outcomes share
$\alpha=\sqrt n\,m^{(1-n)/2}$.  The two probes therefore read the real and
imaginary parts of every selected complex functional.

For completeness, the underlying block inversion is especially simple.
For a non-antipodal separation, set $z_a=H_{a,a-c}$.  Equation
\eqref{eq:general-dark-functional} becomes
\begin{equation}
 \frac{L_{p,c}(H)}{\alpha}=(n-1)z_p+z_{p+c}.
\label{eq:cyclic-block}
\end{equation}
The map is $((n-1)I+P_c)z$ with $P_c$ a permutation.  A vector in its
kernel would obey $(n-1)\|z\|=\|z\|$, and hence vanishes because $n>2$.
For the antipodal class, with $z_p=x_p+\ii y_p$,
\begin{equation}
 \frac{L_{p,c}(H)}{\alpha}
 =(n-1)z_p+z_p^*=nx_p+\ii(n-2)y_p,
\label{eq:antipodal-block}
\end{equation}
which is also invertible.  The cyclic edge classes partition all
off-diagonal Hermitian entries, proving the theorem.  A detailed derivation
is in Appendix~\ref{app:cat-certificate}.

One dark event supplies at most two real rows by Corollary~1.  Consequently
the theorem's $m(m-1)/2$ events and $m(m-1)$ scalar contrasts attain the
corresponding lower bounds \emph{within this weak-probe regular-local
framework}.  This does not prove global injectivity, nor does it prove that
four programmed settings are necessary among more general protocols.  After
division by the known reference magnitudes in
Eq.~\eqref{eq:general-reference-values}, each cyclic block has condition
number at most $n/(n-2)$.  Without that row calibration, the worst
antipodal gain ratio is at most $[n/(n-2)]^2$.

For $m\geq3$, the smallest member of the family is $n=m$; for $m=2$ the
first allowed choice is $n=4$.  At $n=m$, the common squared amplitude
normalization is
\begin{equation}
 \alpha^2=m^{2-m}.
\label{eq:alpha-scaling}
\end{equation}
This exponential-in-$m$ factor is a probability-rate cost, not a rank or
conditioning failure.  With the dense unnormalized probes above, the
probe-only dark-event probability scales as
$p_{\rm ref}=O(\epsilon^2m^{4-m})$ at $n=m$; if the probe operators are
normalized to fixed spectral size it scales instead as
$O(\epsilon^2m^{2-m})$.  These estimates caution against interpreting the
algebraic all-mode theorem as an immediate large-$m$ sample-complexity
advantage.

\subsection{A sparse exact $F_4$ implementation}

Set $m=n=4$, $\ell=0$:
\begin{equation}
 |\Psi_{\rm cat}\rangle
 =\frac{|4000\rangle+|0400\rangle+|0040\rangle+|0004\rangle}{2}.
\label{eq:f4-cat}
\end{equation}
Equation~\eqref{eq:cat-sector} makes the 25 outputs with
$Q(\bm s)\neq0$ exactly dark.  Use the two probe generators
\begin{equation}
 H_X=X_{01}-X_{02},\qquad
 H_Y=Y_{01}-Y_{02}.
\label{eq:cat-probes}
\end{equation}
They are the two quadratures coupling mode $0$ to the antisymmetric
supermode $(|1\rangle-|2\rangle)/\sqrt2$.

In the column orders
\begin{align}
 \btheta_X&=(x_{01},x_{02},x_{03},x_{12},x_{13},x_{23}),\nonumber\\
 \btheta_Y&=(y_{01},y_{02},y_{03},y_{12},y_{13},y_{23}),
\end{align}
six $H_X$ contrasts and six $H_Y$ contrasts have the limiting matrices
\begin{equation}
 J_X=
 \begin{pmatrix}
0&-\frac32&0&0&-\frac32&0\\
\frac34&0&0&\frac34&0&\frac32\\
\frac94&0&0&\frac34&0&0\\
\frac34&0&\frac34&0&0&\frac32\\
0&-4&0&0&0&0\\
\frac34&0&\frac32&\frac34&0&0
\end{pmatrix},
\label{eq:jx}
\end{equation}
and
\begin{equation}
 J_Y=
 \begin{pmatrix}
0&-\frac32&0&0&-\frac32&0\\
\frac34&0&0&\frac34&0&\frac32\\
0&-\frac32&0&0&\frac32&0\\
\frac94&0&0&\frac34&0&0\\
\frac34&0&-\frac34&0&0&\frac32\\
\frac34&0&-\frac32&\frac34&0&0
\end{pmatrix}.
\label{eq:jy}
\end{equation}
The $H_X$ output order is
\begin{equation}
(0022),(0112),(0310),(1021),(1030),(1102),
\label{eq:x-outputs}
\end{equation}
and the $H_Y$ output order is
\begin{equation}
(0022),(0112),(0220),(0310),(1021),(1102).
\label{eq:y-outputs}
\end{equation}
Every omitted cross block is exactly zero.  The exact determinants are
\begin{align}
 \det J_X&=-\frac{243}{8},&
 \det J_Y&=\frac{729}{32},\nonumber\\
 \det\operatorname{diag}(J_X,J_Y)
 &=-\frac{177147}{256}.
\label{eq:cat-determinants}
\end{align}
The Gaussian-integer derivation is in
Appendix~\ref{app:cat-certificate}.

This sparse four-mode certificate is an implementation alternative to
Theorem~3.  It uses seven distinct outcomes rather than the six in the
general construction, while retaining twelve scalar responses and four
programmed settings because all outcomes for a probe are recorded in
parallel: $\pm\epsilon H_X$ and $\pm\epsilon H_Y$.  The raw
limiting condition number is $7.05$; row normalization gives $4.86$, and
ideal background-free Poisson whitening gives $4.79$.  These are local
design diagnostics, not hardware-independent resource guarantees.

\section{Finite probe angle and count statistics}
\label{sec:finite}

The matrices \eqref{eq:jx}--\eqref{eq:jy} are the
$\epsilon\rightarrow0$ Jacobian of $R$, not exact finite-angle
Jacobians.  For the cat input, exact finite-angle probabilities are
available from
\begin{equation}
 A_{\bm s}(U)
 =\frac12\sqrt{\frac{4!}{\prod_ks_k!}}
 \sum_{j=0}^{3}\prod_{k=0}^{3}U_{kj}^{s_k},
 \qquad p_{\bm s}(U)=|A_{\bm s}(U)|^2 .
\label{eq:cat-finite-amplitude}
\end{equation}
For the probes \eqref{eq:cat-probes}, $H^3=2H$, so
\begin{equation}
 \ee^{\ii\epsilon H}
 =I+\frac{\ii\sin(\sqrt2\epsilon)}{\sqrt2}H
 +\frac{\cos(\sqrt2\epsilon)-1}{2}H^2 .
\label{eq:probe-exponential}
\end{equation}

Appendix~\ref{app:signed-proof} proves that the central-contrast Jacobian
has no $O(\epsilon)$ error: its first correction is $O(\epsilon^2)$.
For the selected design, exact finite-angle evaluation gives the compact
diagnostic in Table~\ref{tab:finite-angle}.
\begin{table}[t]
\caption{Finite-angle Jacobian bias and spectral condition number in the
ideal coherent model.  Bias is relative Frobenius error from the limiting
matrix.}
\label{tab:finite-angle}
\begin{ruledtabular}
\begin{tabular}{ccc}
$\epsilon$ & relative bias & $\kappa_2(J_\epsilon)$\\
\hline
0.010 & $7.74\times10^{-4}$ & 7.043\\
0.020 & $3.09\times10^{-3}$ & 7.033\\
0.050 & $1.92\times10^{-2}$ & 6.968\\
0.100 & $7.46\times10^{-2}$ & 6.759
\end{tabular}
\end{ruledtabular}
\end{table}
The favorable condition number at larger angle does not remove the growing
model bias.  Probe-only baseline subtraction in Eq.~\eqref{eq:response} is
also essential: for the selected $H_Y$ outcomes, the largest undivided
baseline begins as $9\epsilon^3/8+O(\epsilon^5)$.

For independent selected Poisson channels,
\begin{equation}
 Y_{a\bm s}\sim\operatorname{Pois}
 \{N_a[p_{a\bm s}(\btheta)+b_{\bm s}]\},
\end{equation}
the Fisher information is
\begin{equation}
 I_{\mu\nu}^{\rm P}
 =\sum_{a,\bm s}N_a
 \frac{\partial_\mu p_{a\bm s}\,\partial_\nu p_{a\bm s}}
 {p_{a\bm s}+b_{\bm s}}.
\label{eq:poisson-fi}
\end{equation}
At a displaced dark event,
$p=c\epsilon^2+O(\epsilon^3)$ and
$\partial_\mu p=d_\mu\epsilon+O(\epsilon^2)$.  The information therefore
has a finite nonzero limit when $b=0$, but is suppressed as
$\epsilon^2/(c\epsilon^2+b)$ above a fixed background floor.  A multinomial
model must additionally include the negative covariance between outcomes
recorded in one setting.  Under any stated likelihood, a Cramér--Rao bound
is conditional on local unbiasedness and does not by itself give a hardware
cost forecast.

This section is an ideal-model statistical layer.  Its formulas and tables
are reproduced by the accompanying finite-angle script and tests.  Loss,
partial distinguishability, source contamination, detector crosstalk,
calibration drift, and optimal shot allocation require an augmented
likelihood before experimental sample complexity can be claimed.

\subsection{Cat-phase confounding}

The coherent input phase is not a benign visibility parameter; it is the
resource that removes the gauge.  If the components of
Eq.~\eqref{eq:f4-cat} acquire unknown phases $\alpha_j$, infinitesimally
they are generated by an input diagonal $D$, and
\begin{equation}
 F_4\ee^{\ii D}
 =\ee^{\ii F_4DF_4^\dagger}F_4.
\label{eq:cat-phase-confounding}
\end{equation}
The three relative cat phases are therefore exactly confounded with the
three device directions in Eq.~\eqref{eq:gauge-basis}.  If those preparation
phases are free nuisance parameters, the augmented model does not separately
identify all twelve device coordinates.  Independent cat-phase calibration
or additional reference states are necessary.  Complete dephasing returns
the experiment to the number-state gauge.  Even with calibrated phases,
the construction is not globally injective: for a diagonal
$D=\operatorname{diag}(d_0,\ldots,d_{m-1})$, the discrete conditions
\begin{equation}
 n(d_j-d_0)\in2\pi\mathbb Z\qquad(1\leq j<m)
\label{eq:cat-global-alias}
\end{equation}
leave the equal-phase $n$-photon cat unchanged up to a global phase.
Consequently the corresponding finite transported transformations
$F_m\exp(\ii D)F_m^\dagger$ are exact global aliases.  They are discrete
near a fixed calibrated representative and therefore do not contradict the
local full-rank theorem.

\section{Internal coherent errors}
\label{sec:internal}

Let a nominal circuit factor as
\begin{equation}
 U_0=U_LU_{L-1}\cdots U_1.
\end{equation}
If a small coherent generator $H_\ell$ is inserted after layer $\ell$, set
$W_\ell=U_L\cdots U_{\ell+1}$.  Then
\begin{equation}
 U_L\cdots U_{\ell+1}
 \ee^{\ii\delta_\ell H_\ell}
 U_\ell\cdots U_1
 =
 \ee^{\ii\delta_\ell W_\ell H_\ell W_\ell^\dagger}U_0.
\label{eq:single-internal}
\end{equation}
For simultaneous small errors, the first-order output-equivalent generator
is
\begin{equation}
 G_{\rm eff}
 =\sum_\ell\delta_\ell
 W_\ell H_\ell W_\ell^\dagger+O(\|\bm\delta\|^2).
\label{eq:effective-generator}
\end{equation}
Thus the signed-probe reconstruction applies to arbitrary internal coherent,
lossless errors after adjoint propagation.

This does not automatically locate a faulty component.  Different layer
errors can have the same output-equivalent generator, and output-diagonal
projections remain outside our $m(m-1)$-coordinate target.  Layer localization
requires the circuit Jacobian across additional configurations, addressable
controls, or structural priors.  Appendix~\ref{app:internal-proof} gives the
short proof and the multi-error expansion.

\section{Relation to prior work and limitations}
\label{sec:limitations}

The deliberate displacement of a null measurement is a general estimation
principle, not one introduced here.  Girotti \emph{et al.} established its
role in resolving null-measurement non-identifiability and in asymptotically
optimal pure-state estimation \cite{Girotti2024}.  Our contribution is an
explicit multiphoton-count implementation, the exact number-state gauge
theorem, and explicit dimension-saturating local frames for a Fourier
device.

Likewise, Eq.~\eqref{eq:cat-sector} is a known Fourier-cat/cyclic selection
effect \cite{VourdasDunningham2005,Dittel2018PRA}.  The new role assigned to
the cat is operational: it supplies the missing relative phase reference.
Standard Fourier suppression experiments
\cite{Tichy2010,Crespi2016,Dittel2018PRL} and programmable photonic
processors \cite{Arrazola2021,Wang2023} provide the experimental context,
but preparing a phase-stable multimode multiphoton cat is substantially
harder than preparing a number state.  The rank-nine $F_4$ Fock design is
the nearer-term protocol; the cat family is a complete-identifiability
benchmark.

The closest comparator we found for general near-identity unitary
estimation reconstructs all $d^2-1$ parameters and can attain the quantum
Cramér--Rao bound, but uses a target qudit, two control qudits, and
controlled shift and phase operations \cite{EscandonMonardes2024}.  It does
not provide a passive-bosonic Fourier construction from photon-counted dark
outcomes.  To our knowledge, prior work has not combined Fourier-cat
suppression with displaced-null readout to give an explicit all-mode,
dimension-saturating local frame for coherent off-diagonal errors of a
passive multiport.  This is a scoped literature assessment, not a claim to
the first arbitrary-unitary tomography protocol, the first displaced-null
measurement, or a new Fourier-cat suppression law.

The rigorous scope is:
\begin{itemize}
\item local reconstruction around a calibrated nominal interferometer, not
global injectivity or global unitary tomography;
\item $m(m-1)$ off-diagonal coordinates only, not the $m-1$
output-diagonal $SU(m)$ directions;
\item coherent lossless errors with calibrated probes and input phases;
\item two probe directions and four signed settings suffice, but neither a
one-probe nor a fewer-setting impossibility theorem has been proved;
\item the all-mode outcome and scalar-contrast counts, and the $F_4$ Fock
counts, are minimal only within the declared weak-probe regular-local
framework; the finite-angle experiment has not been globally optimized;
\item identification is local and calibrated: global aliases and
off-diagonal device errors confounded with uncalibrated input phases or
probe errors are not excluded;
\item at $n=m$, $\alpha^2=m^{2-m}$, so event rates decay rapidly even
though the row-normalized condition bound is at most $m/(m-2)$;
\item the ideal likelihood does not absorb loss or distinguishability into
a universal ``background'' parameter
\cite{Shchesnovich2015,DufourBuchleitner2026}.
\end{itemize}

\section{Discussion}
\label{sec:discussion}

A multiphoton dark event is not by itself a signed sensor.  The known probe
in Eq.~\eqref{eq:response} supplies a reference amplitude and turns photon
counts into linear local observables.  Which coordinates can then be
reconstructed depends on both the collection of dark amplitude gradients
and the phase reference carried by the input state.

This viewpoint exposes a sharp resource separation.  Number states have an
exact $(m-1)$-dimensional blind subspace at $F_m$; the $F_4$ certificate
saturates the resulting rank-nine bound.  A calibrated cat state lifts that
continuous gauge, and the cyclic-edge construction gives an invertible
$m(m-1)$-coordinate design for every mode number.  Unknown cat phases
recreate precisely the obstruction the state was introduced to remove.  The
distinction is therefore one of identifiability, not simply better
visibility.

The most immediate extensions are a joint SPAM--device likelihood,
lower-cost coherent reference states, circuit-specific internal-error
dictionaries, and sparse or bounded-norm probes that improve large-$m$
count rates.  Any such extension must preserve the separation between local
full rank, finite-angle bias, finite-shot efficiency, and global uniqueness.

\begin{acknowledgments}
The author used OpenAI Codex as a research and writing assistant for
symbolic exploration, exact-code generation, literature discovery, and
editorial revision.  All claims and calculations remain the author's
responsibility.
\end{acknowledgments}

\appendix

\section{Proof of the signed-probe and finite-angle theorems}
\label{app:signed-proof}

Put $|\phi_e\rangle=\widehat U_0|\psi_e\rangle$ and suppress the event
subscript.  For a fixed probe $\widehat h$, define
\begin{align}
 a(t)&=\langle\bm s|\ee^{\ii t\widehat h}|\phi\rangle,\nonumber\\
 b_\mu(t)&=\langle\bm s|\ee^{\ii t\widehat h}
 \ii\widehat G_\mu|\phi\rangle,\nonumber\\
 g_\mu(t)&=\left.\partial_{\theta_\mu}
 |\langle\bm s|\ee^{\ii t\widehat h}
 \ee^{\ii\widehat G(\btheta)}|\phi\rangle|^2
 \right|_{\btheta=0}\nonumber\\
 &=2\operatorname{Re}\{a(t)^*b_\mu(t)\}.
\label{eq:appendix-g}
\end{align}
Darkness gives $a(0)=0$.  Let
\begin{equation}
 a_r=\langle\bm s|(\ii\widehat h)^r|\phi\rangle,\qquad
 b_{r\mu}=\langle\bm s|(\ii\widehat h)^r
 \ii\widehat G_\mu|\phi\rangle .
\end{equation}
The derivative of the response \eqref{eq:response} is
\begin{equation}
 \left.\partial_{\theta_\mu}R\right|_0
 =\frac{g_\mu(\epsilon)-g_\mu(-\epsilon)}{\epsilon}.
\end{equation}
Taylor expansion of $a$ and $b$ gives
\begin{align}
 \left.\partial_{\theta_\mu}R\right|_0
={}&4\operatorname{Re}(a_1^*b_{0\mu})\nonumber\\
&+4\epsilon^2\operatorname{Re}\left\{
\frac{a_1^*b_{2\mu}}2+\frac{a_2^*b_{1\mu}}2
+\frac{a_3^*b_{0\mu}}6\right\}
+O(\epsilon^4).
\label{eq:finite-expansion}
\end{align}
Since $a_1=\ell\bm h$ and $b_{0\mu}=\ell\be_\mu$,
Eq.~\eqref{eq:signed-probe-theorem} follows.  Writing
$\ell=\bm u+\ii\bm v$ as a complex row, its probe row is a real linear
combination of $\bm u$ and $\bm v$, which proves
Eq.~\eqref{eq:two-row-span}.

For a norm bound, let
\begin{equation}
 M=\|\widehat h\|\leq n\|h\|,\qquad
 L_\mu=\|\widehat G_\mu\|\leq n\|G_\mu\|.
\end{equation}
Leibniz's rule applied to Eq.~\eqref{eq:appendix-g} yields
\begin{equation}
 |g_\mu'''(t)|
 \leq2\sum_{k=0}^{3}\binom3kM^kM^{3-k}L_\mu
 =16M^3L_\mu .
\end{equation}
The central-difference remainder theorem therefore gives, for the
convention \eqref{eq:response},
\begin{equation}
 \left|
 \left.\partial_{\theta_\mu}R\right|_0
 -4\operatorname{Re}(a_1^*b_{0\mu})
 \right|
 \leq\frac{16}{3}\epsilon^2M^3L_\mu .
\label{eq:finite-bound}
\end{equation}
This bound is conservative but uniform.  Retaining the explicit
$\epsilon^2$ term in Eq.~\eqref{eq:finite-expansion} leaves an
$O(\epsilon^4)$ remainder.  Baseline subtraction does not change the
derivative with respect to $\btheta$, but it removes the finite-angle
probe-only offset from the observed response.

\section{Fock gauge proof and exact rank-nine certificate}
\label{app:fock-certificate}

For $D=\operatorname{diag}(d_0,\ldots,d_{m-1})$,
\begin{equation}
 (F_mDF_m^\dagger)_{pp}
 =\frac1m\sum_jd_j=\frac{\Tr D}{m}=0.
\end{equation}
Conjugation commutes with exponentiation, so
$\ee^{\ii K_D}F_m=F_m\ee^{\ii D}$ at the single-particle level and hence
in every symmetric photon sector.  A number state is an eigenvector of
$\widehat D=\sum_jd_j\widehat n_j$, proving
Eq.~\eqref{eq:gauge-no-go}.  Appending a known unitary $V$ preserves the
global phase, so no subsequent number measurement removes the
obstruction.  The map from traceless $D$ to $K_D$ is injective and its
dimension is $m-1$.

For completeness, the $F_4$ finite certificate can be evaluated with
Gaussian integers.  Let
\begin{equation}
 Z_{\bm r,\bm t}=\per(2F_4)[\bm t,\bm r],\qquad
 D_{\bm r,\bm s}=2^n
 \sqrt{\prod_jr_j!\prod_ks_k!}.
\end{equation}
For a dark target, the first amplitude derivatives are
\begin{align}
 \ell X_{pq}
 &=\frac{\ii}{D_{\bm r,\bm s}}
 \left(s_qZ_{\bm r,\bm s+\be_p-\be_q}
 +s_pZ_{\bm r,\bm s-\be_p+\be_q}\right),\nonumber\\
 \ell Y_{pq}
 &=\frac{1}{D_{\bm r,\bm s}}
 \left(s_pZ_{\bm r,\bm s-\be_p+\be_q}
 -s_qZ_{\bm r,\bm s+\be_p-\be_q}\right).
\label{eq:neighboring-permanents}
\end{align}
Negative occupations are omitted.  These follow directly from the Fock
ladder factors; the square roots cancel against the changed-state
normalization.

Substitution into Eq.~\eqref{eq:signed-probe-theorem}, for the row order of
Table~\ref{tab:fock-design}, gives
\begin{widetext}
\begin{equation}
J_F=
\begin{pmatrix}
0&0&0&0&\frac32&0&0&0&0&0&-\frac32&0\\
0&0&0&0&0&\frac32&0&0&0&0&0&\frac32\\
0&0&1&0&0&0&0&0&-1&0&0&0\\
0&0&0&1&0&0&0&0&0&-1&0&0\\
0&0&0&0&0&0&\frac32&0&0&0&-\frac32&0\\
0&0&0&0&0&0&0&\frac32&0&0&0&-\frac32\\
-1&0&0&0&0&0&-1&0&0&0&2&0\\
0&-1&0&0&0&0&0&-1&0&0&0&2\\
0&0&0&1&0&0&0&0&0&1&0&0
\end{pmatrix}.
\label{eq:fock-jacobian}
\end{equation}
\end{widetext}
Each row annihilates the coordinate vectors of $K_A,K_B,K_C$.
The minor in columns
\begin{equation}
(x_{01},y_{01},x_{02},y_{02},x_{03},y_{03},x_{12},y_{12},y_{13})
\end{equation}
has determinant $81/8$.  Hence $\operatorname{rank}J_F=9$ and the three
independent gauge vectors form its complete nullspace.

\section{Fourier-cat frame proof and sparse rank-twelve certificate}
\label{app:cat-certificate}

For the component $|n\be_j\rangle$, every term in the repeated-column
permanent has the same phase.  Equation~\eqref{eq:fock-amplitude} gives
\begin{equation}
 {\cal A}_{|n\be_j\rangle,\bm s}(F_m)
 =\sqrt{\frac{n!}{\prod_ks_k!}}\,
 m^{-n/2}\omega^{jQ(\bm s)} .
\end{equation}
Multiplying by the coefficients in Eq.~\eqref{eq:general-cat} and summing
over $j$ gives
\begin{equation}
 \frac1{\sqrt m}\sum_{j=0}^{m-1}
 \omega^{j[Q(\bm s)-\ell]}
 =\sqrt m\,\delta_{Q(\bm s),\ell},
\end{equation}
which proves Eq.~\eqref{eq:cat-sector}.

We next derive Eq.~\eqref{eq:general-dark-functional} and complete the
all-mode proof.  For a general single-particle matrix $U$, the cat amplitude
is the polynomial
\begin{equation}
 A_{\bm s}(U)=
 \sqrt{\frac{n!}{\prod_as_a!}}\,\frac1{\sqrt m}
 \sum_{j=0}^{m-1}\prod_{a=0}^{m-1}U_{aj}^{s_a}.
\label{eq:general-cat-polynomial}
\end{equation}
Along $U(\theta)=\ee^{\ii\theta H}F_m$,
$\dot U_{aj}=\ii\sum_bH_{ab}F_m(b,j)$.  Differentiating
Eq.~\eqref{eq:general-cat-polynomial} at zero produces terms proportional
to
\begin{equation}
 \sum_{j=0}^{m-1}\omega^{j[Q(\bm s)+b-a]}
 =m\,\delta_{b,a-c}
\end{equation}
when $Q(\bm s)=c$.  Collecting the factorial and Fourier normalizations
gives
\begin{equation}
 \left.\frac{d}{d\theta}A_{\bm s}
   (\ee^{\ii\theta H}F_m)\right|_{\theta=0}
 =\ii\alpha_{\bm s}\sum_as_aH_{a,a-c}
 =\ii L_{\bm s}(H),
\label{eq:general-cat-gradient}
\end{equation}
as asserted.

For the selected outcome \eqref{eq:selected-general-outcomes}, all
factorials are the same and
$\alpha=\sqrt n\,m^{(1-n)/2}$.  In a non-antipodal cyclic class, the
$m$ complex coordinates $z_a=H_{a,a-c}$ are independent, because their
Hermitian conjugates belong to the unselected class $-c$.  The selected
functionals form the complex matrix $(n-1)I+P_c$.  If it annihilates $z$,
unitarity of the permutation gives
\begin{equation}
 (n-1)\|z\|=\|P_cz\|=\|z\|,
\end{equation}
so $z=0$ for $n>2$.  In the self-conjugate antipodal class, Hermiticity
instead reduces each edge to $z_p=x_p+\ii y_p$, and its functional is
$nx_p+\ii(n-2)y_p$.  This is invertible over the reals.  Each undirected
off-diagonal edge appears in exactly one selected cyclic class, so these
blocks are a direct-sum decomposition of ${\cal H}_{\rm off}(m)$.

It remains to show that the complex functionals are observable from counts.
For the probes defined in Eq.~\eqref{eq:oriented-probe}, direct substitution
gives Eq.~\eqref{eq:general-reference-values}.  Since
$\ell=\ii L$, Theorem~1 reduces here to
\begin{equation}
 D R_{\bm s,h}(0)[H]
 =4\operatorname{Re}\{L_{\bm s}(h)^*L_{\bm s}(H)\}.
\label{eq:general-cat-response}
\end{equation}
The $R$ and $T$ rows are therefore nonzero multiples of the real and
imaginary parts of every $L_{\bm s}$.  The preceding direct-sum argument
proves full rank.  There are exactly $m(m-1)$ such real rows.  Since
Corollary~1 bounds one event by two rows, the event and scalar-response
counts meet the weak-probe regular-local lower bounds.  Analytic dependence
on $\epsilon$ and persistence of a nonzero determinant then establish the
finite-angle local statement through the ordinary inverse-function theorem.

The conditioning claim follows from the singular values
$|n-1+\lambda|$, $|\lambda|=1$, of each non-antipodal block, and from the
two gains $n,n-2$ in an antipodal block.  Their ratio is at most
$n/(n-2)$ after known reference-row normalization.  Finally,
$\alpha^2=n m^{1-n}$, which becomes Eq.~\eqref{eq:alpha-scaling} when
$n=m$.

We now certify the sparse $F_4$ matrices without numerical zero tests.  For any
four-mode unitary,
\begin{equation}
 A_{\bm s}(U)
 =\frac12\sqrt{\frac{4!}{\prod_ks_k!}}
 \sum_{j=0}^{3}\prod_{k=0}^{3}U_{kj}^{s_k}.
\end{equation}
Its derivative in the direction $\dot U_\mu=\ii G_\mu F_4$ is
\begin{align}
 \dot A_{\bm s,\mu}
 ={}&\frac12\sqrt{\frac{4!}{\prod_ks_k!}}
 \sum_{j=0}^{3}\sum_{k:s_k>0}
 s_k(\dot U_\mu)_{kj}U_{kj}^{s_k-1}
 \prod_{\substack{r=0\\r\neq k}}^3U_{rj}^{s_r}.
\label{eq:cat-gradient}
\end{align}
At $U=F_4$, all entries lie in $\mathbb Q(\ii)$ times the displayed
normalization.  The limiting physical response is
\begin{equation}
 J_{\bm s,\mu}
 =4\operatorname{Re}
 \left(\dot A_{\bm s,h}^*\dot A_{\bm s,\mu}\right).
\label{eq:cat-response-entry}
\end{equation}
Evaluating Eq.~\eqref{eq:cat-gradient} for the outputs
\eqref{eq:x-outputs}--\eqref{eq:y-outputs} gives
Eqs.~\eqref{eq:jx}--\eqref{eq:jy}.  Direct expansion gives
\begin{equation}
 \det J_X=-243/8,\qquad \det J_Y=729/32.
\end{equation}
The $H_X$ rows have zero entries in all six $Y$ columns and the $H_Y$
rows have zero entries in all six $X$ columns.  Thus the full determinant
is their product, Eq.~\eqref{eq:cat-determinants}.  This is an exact
rational certificate.  The released verification independently constructs
the same rows from Gaussian-integer permanent moments and asserts every
cross-block zero.

\section{Internal-error propagation}
\label{app:internal-proof}

Write $U_0=W_\ell V_\ell$, with
$W_\ell=U_L\cdots U_{\ell+1}$ and
$V_\ell=U_\ell\cdots U_1$.  Then
\begin{align}
 W_\ell\ee^{\ii\delta H_\ell}V_\ell
 &=
 W_\ell\ee^{\ii\delta H_\ell}W_\ell^\dagger U_0\nonumber\\
 &=\ee^{\ii\delta W_\ell H_\ell W_\ell^\dagger}U_0,
\end{align}
which proves Eq.~\eqref{eq:single-internal} exactly for one inserted error.
For several errors, expanding each exponential to first order and collecting
terms gives Eq.~\eqref{eq:effective-generator}.  Commutators enter at second
order, so the layer-to-output dictionary is linear only in the local regime.

\bibliography{references}

\end{document}
